\section{Relationship to the preserved manuscript}\label{app:checkpoint}
The previous manuscript, \emph{When Kernel Affine Hull Geometry Helps
Prediction: Local Coordinates, Decoder Capacity, and Conditional Gains},
is preserved with its reproducibility archive at commit
\texttt{a67c40190458bab7b3e56b4fac76770798d074b3}. The new extension records
SHA-256 hashes for all 217 previously tracked repository files and the seven
submission artifacts. The original study is neither relabeled as natural-data
confirmation nor pooled with the new source families. Its synthetic
experiments provide motivation and mechanism controls; the new confirmation
is based exclusively on the reserved natural sources.

At the 256-landmark setting, the earlier Van der Pol experiment reported mean
MSEs of 0.00165415 for raw inputs, 0.00107800 for affine local coordinates, and
0.00087197 for KAHM local coordinates. Across its separately defined capacity
contrast family, four comparisons supported practical benefit, six supported
harm, and two were unresolved. These are archived results, not newly estimated
effects or additional independent forecasting domains. The previous Airfoil
experiment was static regression and is not included as a natural forecasting
task in the present study.

The new implementation preserves the local KAHM definition while making
standalone alternative fitting and alpha-path reuse explicit. Its stable
Gaussian query calculation and raw-augmentation choices are specified in
Section~\ref{sec:nf-methods}. The former manuscript remains byte-identical;
the first-page disclosure in this new draft also addresses the placement
issue identified in the previous submission-format review.

\section{Reconstruction of the experiment}\label{app:reproduction}
The repository entry point is
\texttt{extensions/natural\_forecasting}. From the repository root, install
the extension's pinned requirements in Python 3.12, then run the synthetic
unit tests. The data command verifies the eight checksummed archive snapshots.
The experiment is separated into preparation, initial freeze, development,
diagnostic freeze, confirmation, and analysis:
\begin{verbatim}
python -m extensions.natural_forecasting.run prepare
python -m extensions.natural_forecasting.run initial-freeze
python -m extensions.natural_forecasting.run run --role development
python -m extensions.natural_forecasting.run freeze-confirmation
python -m extensions.natural_forecasting.run run --role confirmation
python -m extensions.natural_forecasting.run analyze
\end{verbatim}
Each command accepts a separate run directory. A new reproduction makes its
own freeze; the supplied reference freeze is immutable. Confirmation cannot
start without a matching code, dependency, roster, and diagnostic state.
Interrupted complete tasks can be resumed only when their prediction artifacts
match their recorded hashes. A failed task is retained and stops the run;
silently replacing a task or dropping a failed method is not permitted by the
analysis. The supplementary archive includes the exact reference results and
the metadata needed to audit them without refitting every predictor.

The two random seeds are 17 and 43. NumPy, SciPy, scikit-learn, statsmodels,
pandas, and threadpoolctl versions, thread settings, visible processor
information, and the separate library initialization cost are recorded in
the freeze. These details support reproducibility of the computation.
Wall-clock measurements remain specific to the declared shared CPU
environment, and process high-water memory is not an incremental allocation
measurement.

\begin{table}[t]
\centering\small\setlength{\tabcolsep}{4pt}
\caption{Completed experiment audit. Counts of forecasts, tasks and tests serve reproducibility checks; they do not increase the number of independent source families. The report builder reads existing predictions and validation evidence without rerunning forecasting or diagnostic inference.}\label{tab:validation-audit}
\begin{tabular}{lr}
\toprule
Audit item & Recorded result \\
\midrule
Final runner validation & passed \\
Development original source families & 4 \\
Confirmation original source families & 4 \\
Development / confirmation task runs & 128 / 80 \\
Total task runs & 208 \\
Saved candidate forecast arrays / dimensions checked & 2496 \\
Candidate losses independently recomputed by runner & 2496 \\
Targets and origins equal frozen source windows & True \\
Historical tracked files verified unchanged & 217 \\
Scientific protocol and both stage freezes intact & True \\
Unit tests in supplied evidence & 53 \\
\bottomrule
\end{tabular}
\end{table}


\section{Additional source and capacity results}\label{app:full-results}
All declared comparisons are retained below. The two context--horizon pairs
are coupled settings, so the tables do not separately identify a causal
context-length or horizon effect. Likewise, the two capacity settings change
both the cap on training windows and specified model sizes. Their timings are
observed costs rather than enforced equal-time budgets.
{\small\setlength{\tabcolsep}{4pt}
\begin{longtable}{lrrrrr}
\caption{Development scaled MSE by original source family. The final column gives equal-source aggregation; source columns retain the declared within-source weights. Development diagnostics use leave-one-source-family-out fits; confirmation diagnostics use the frozen development fit. Always-policy rows duplicate their corresponding standalone methods deliberately.}\label{tab:source-development-loss}\\
\toprule
Method or policy & FRED-MD & NN5 & Tourism & Electricity & Equal-source \\
\midrule
\endfirsthead
\toprule
Method or policy & FRED-MD & NN5 & Tourism & Electricity & Equal-source \\
\midrule
\endhead
Raw ridge & 0.5601 & 0.8584 & 2.865 & 0.4243 & 1.177 \\
Raw RBF & 10.46 & 0.7552 & 3.476 & 0.4388 & 3.783 \\
Affine local & 2.548 & 0.8795 & 3.769 & 0.4535 & 1.912 \\
Locality local & 23.83 & 0.8148 & 3.778 & 0.4112 & 7.209 \\
KAHM local & 12.23 & 0.9163 & 3.636 & 0.4525 & 4.309 \\
Raw + affine & 1.409 & 0.7614 & 3.172 & 0.4481 & 1.448 \\
Raw + locality & 10.18 & 0.7474 & 3.046 & 0.4418 & 3.603 \\
Raw + KAHM & 4.254 & 0.7704 & 3.229 & 0.4539 & 2.177 \\
Seasonal naive & 2.269 & 1.089 & 2.982 & 0.5507 & 1.723 \\
Holt & 0.2409 & 1.358 & 8.228 & 0.6984 & 2.631 \\
Histogram boosting & 34.02 & 0.8032 & 4.786 & 0.4404 & 10.01 \\
MLP & 1.022 & 0.8475 & 2.77 & 0.4715 & 1.278 \\
Always raw & 10.46 & 0.7552 & 3.476 & 0.4388 & 3.783 \\
Always raw + KAHM & 4.254 & 0.7704 & 3.229 & 0.4539 & 2.177 \\
Validation: matched pair & 5.267 & 0.7335 & 3.229 & 0.4415 & 2.418 \\
Validation: raw portfolio & 0.2409 & 0.8021 & 3.014 & 0.4528 & 1.127 \\
Validation: KAHM family & 4.104 & 0.7792 & 3.223 & 0.4572 & 2.141 \\
Validation: full portfolio & 0.2409 & 0.8052 & 3.092 & 0.4675 & 1.152 \\
Cheap probe & 5.106 & 0.742 & 3.25 & 0.4217 & 2.38 \\
Metadata & 4.254 & 0.7704 & 3.229 & 0.4539 & 2.177 \\
Metadata + geometry & 4.254 & 0.7704 & 3.394 & 0.4539 & 2.218 \\
Metadata + probe & 4.157 & 0.759 & 3.338 & 0.4342 & 2.172 \\
Metadata + probe + geometry & 4.215 & 0.7704 & 3.338 & 0.4342 & 2.189 \\
\bottomrule
\end{longtable}
}

{\small\setlength{\tabcolsep}{4pt}
\begin{longtable}{lrrrrr}
\caption{Development warm-library wall seconds by original source family. The final column gives equal-source aggregation; source columns retain the declared within-source weights. Development diagnostics use leave-one-source-family-out fits; confirmation diagnostics use the frozen development fit. Always-policy rows duplicate their corresponding standalone methods deliberately.}\label{tab:source-development-cost}\\
\toprule
Method or policy & FRED-MD & NN5 & Tourism & Electricity & Equal-source \\
\midrule
\endfirsthead
\toprule
Method or policy & FRED-MD & NN5 & Tourism & Electricity & Equal-source \\
\midrule
\endhead
Raw ridge & 0.3731 & 0.1631 & 0.1888 & 0.1873 & 0.2281 \\
Raw RBF & 0.1929 & 0.1912 & 0.1834 & 0.1692 & 0.1842 \\
Affine local & 0.3285 & 0.2805 & 0.2983 & 0.2863 & 0.2984 \\
Locality local & 0.5431 & 0.29 & 0.2784 & 0.2649 & 0.3441 \\
KAHM local & 0.3705 & 0.3696 & 0.3735 & 0.3369 & 0.3626 \\
Raw + affine & 0.2923 & 0.2998 & 0.2979 & 0.288 & 0.2945 \\
Raw + locality & 0.3185 & 0.2967 & 0.2988 & 0.2533 & 0.2918 \\
Raw + KAHM & 0.4275 & 0.3807 & 0.3252 & 0.3024 & 0.3589 \\
Seasonal naive & 0.2125 & 0.177 & 0.1834 & 0.1878 & 0.1902 \\
Holt & 0.3224 & 0.42 & 0.3682 & 0.316 & 0.3567 \\
Histogram boosting & 3.086 & 4.954 & 1.729 & 1.819 & 2.897 \\
MLP & 1.278 & 0.9551 & 0.8521 & 0.5508 & 0.9089 \\
Always raw & 0.1929 & 0.1912 & 0.1834 & 0.1692 & 0.1842 \\
Always raw + KAHM & 0.4275 & 0.3807 & 0.3252 & 0.3024 & 0.3589 \\
Validation: matched pair & 0.5562 & 0.5313 & 0.4929 & 0.4301 & 0.5026 \\
Validation: raw portfolio & 4.066 & 5.306 & 2.641 & 2.719 & 3.683 \\
Validation: KAHM family & 0.692 & 0.6568 & 0.6115 & 0.5511 & 0.6278 \\
Validation: full portfolio & 5.861 & 6.513 & 4.138 & 4.075 & 5.147 \\
Cheap probe & 0.4739 & 0.4615 & 0.4004 & 0.3328 & 0.4172 \\
Metadata & 0.4284 & 0.3815 & 0.3258 & 0.3031 & 0.3597 \\
Metadata + geometry & 0.5195 & 0.465 & 0.3474 & 0.3744 & 0.4266 \\
Metadata + probe & 0.5029 & 0.474 & 0.2891 & 0.3411 & 0.4018 \\
Metadata + probe + geometry & 0.5284 & 0.4828 & 0.2891 & 0.341 & 0.4103 \\
\bottomrule
\end{longtable}
}

{\small\setlength{\tabcolsep}{4pt}
\begin{longtable}{lrrrrr}
\caption{Confirmation scaled MSE by original source family. The final column gives equal-source aggregation; source columns retain the declared within-source weights. Development diagnostics use leave-one-source-family-out fits; confirmation diagnostics use the frozen development fit. Always-policy rows duplicate their corresponding standalone methods deliberately.}\label{tab:source-confirmation-loss}\\
\toprule
Method or policy & Traffic & Hospital & Saugeen & US births & Equal-source \\
\midrule
\endfirsthead
\toprule
Method or policy & Traffic & Hospital & Saugeen & US births & Equal-source \\
\midrule
\endhead
Raw ridge & 0.5737 & 0.8461 & 0.522 & 1.177 & 0.7798 \\
Raw RBF & 0.5739 & 0.8228 & 0.6106 & 1.223 & 0.8075 \\
Affine local & 0.5796 & 0.928 & 0.6462 & 1.626 & 0.9449 \\
Locality local & 0.5719 & 0.7702 & 0.6504 & 1.593 & 0.8964 \\
KAHM local & 0.5775 & 0.9029 & 0.6652 & 1.596 & 0.9354 \\
Raw + affine & 0.5764 & 0.9198 & 0.6068 & 1.28 & 0.8458 \\
Raw + locality & 0.572 & 0.9356 & 0.6002 & 1.277 & 0.8462 \\
Raw + KAHM & 0.5768 & 0.9238 & 0.6012 & 1.266 & 0.8419 \\
Seasonal naive & 0.6484 & 0.7848 & 1.023 & 1.743 & 1.05 \\
Holt & 0.6842 & 1.003 & 0.6079 & 7.197 & 2.373 \\
Histogram boosting & 0.6524 & 1.025 & 0.5386 & 1.606 & 0.9555 \\
MLP & 0.6811 & 1.15 & 0.5365 & 1.147 & 0.8785 \\
Always raw & 0.5739 & 0.8228 & 0.6106 & 1.223 & 0.8075 \\
Always raw + KAHM & 0.5768 & 0.9238 & 0.6012 & 1.266 & 0.8419 \\
Validation: matched pair & 0.5738 & 0.8236 & 0.6042 & 1.237 & 0.8097 \\
Validation: raw portfolio & 0.6087 & 1.116 & 0.5639 & 1.167 & 0.8639 \\
Validation: KAHM family & 0.577 & 0.9182 & 0.6012 & 1.266 & 0.8406 \\
Validation: full portfolio & 0.6122 & 1.099 & 0.5695 & 1.225 & 0.8764 \\
Cheap probe & 0.5761 & 0.8211 & 0.6253 & 1.235 & 0.8143 \\
Metadata & 0.5739 & 0.8228 & 0.6012 & 1.266 & 0.816 \\
Metadata + geometry & 0.5739 & 0.8228 & 0.6012 & 1.266 & 0.816 \\
Metadata + probe & 0.5761 & 0.8189 & 0.6106 & 1.223 & 0.8071 \\
Metadata + probe + geometry & 0.5761 & 0.8189 & 0.6106 & 1.223 & 0.8071 \\
\bottomrule
\end{longtable}
}

{\small\setlength{\tabcolsep}{4pt}
\begin{longtable}{lrrrrr}
\caption{Confirmation warm-library wall seconds by original source family. The final column gives equal-source aggregation; source columns retain the declared within-source weights. Development diagnostics use leave-one-source-family-out fits; confirmation diagnostics use the frozen development fit. Always-policy rows duplicate their corresponding standalone methods deliberately.}\label{tab:source-confirmation-cost}\\
\toprule
Method or policy & Traffic & Hospital & Saugeen & US births & Equal-source \\
\midrule
\endfirsthead
\toprule
Method or policy & Traffic & Hospital & Saugeen & US births & Equal-source \\
\midrule
\endhead
Raw ridge & 0.1454 & 0.1959 & 0.1905 & 0.2204 & 0.188 \\
Raw RBF & 0.1644 & 0.1636 & 0.1693 & 0.1779 & 0.1688 \\
Affine local & 0.2424 & 0.2732 & 0.4038 & 1.082 & 0.5002 \\
Locality local & 0.2311 & 0.254 & 0.4021 & 0.2721 & 0.2898 \\
KAHM local & 0.2579 & 0.2764 & 0.3496 & 0.4035 & 0.3218 \\
Raw + affine & 0.3969 & 0.2809 & 0.354 & 0.3199 & 0.3379 \\
Raw + locality & 0.2223 & 0.2752 & 0.3166 & 0.3022 & 0.2791 \\
Raw + KAHM & 0.2674 & 0.4235 & 0.4676 & 0.3911 & 0.3874 \\
Seasonal naive & 0.1551 & 0.164 & 0.1874 & 0.1393 & 0.1614 \\
Holt & 0.2759 & 0.3535 & 0.4118 & 0.3386 & 0.3449 \\
Histogram boosting & 1.774 & 1.39 & 17.59 & 14.89 & 8.91 \\
MLP & 0.4046 & 0.392 & 2.264 & 2.79 & 1.463 \\
Always raw & 0.1644 & 0.1636 & 0.1693 & 0.1779 & 0.1688 \\
Always raw + KAHM & 0.2674 & 0.4235 & 0.4676 & 0.3911 & 0.3874 \\
Validation: matched pair & 0.3909 & 0.5408 & 0.5273 & 0.4656 & 0.4812 \\
Validation: raw portfolio & 2.276 & 2.19 & 15.94 & 12.95 & 8.338 \\
Validation: KAHM family & 0.4622 & 0.6403 & 0.7116 & 0.6829 & 0.6242 \\
Validation: full portfolio & 3.595 & 3.661 & 17.63 & 14.9 & 9.948 \\
Cheap probe & 0.3029 & 0.2768 & 0.4534 & 0.3531 & 0.3465 \\
Metadata & 0.1651 & 0.1644 & 0.4686 & 0.3923 & 0.2976 \\
Metadata + geometry & 0.2216 & 0.2212 & 0.5797 & 0.4555 & 0.3695 \\
Metadata + probe & 0.2958 & 0.457 & 0.3033 & 0.2673 & 0.3309 \\
Metadata + probe + geometry & 0.2958 & 0.457 & 0.3033 & 0.2673 & 0.3309 \\
\bottomrule
\end{longtable}
}

{\small\setlength{\tabcolsep}{4pt}
\begin{longtable}{lrrrr}
\caption{Development capacity comparison. Each capacity changes specified training-window and model-size caps; it is not an equal-time budget. Values are equal-source weighted and retain all prescribed context--horizon settings and seeds.}\label{tab:capacity-development}\\
\toprule
Method or policy & medium MSE & medium wall (s) & small MSE & small wall (s) \\
\midrule
\endfirsthead
\toprule
Method or policy & medium MSE & medium wall (s) & small MSE & small wall (s) \\
\midrule
\endhead
Raw ridge & 1.177 & 0.28 & 1.177 & 0.1761 \\
Raw RBF & 4.582 & 0.1846 & 2.984 & 0.1838 \\
Affine local & 1.989 & 0.3292 & 1.836 & 0.2676 \\
Locality local & 8.238 & 0.4065 & 6.18 & 0.2817 \\
KAHM local & 3.855 & 0.3779 & 4.763 & 0.3473 \\
Raw + affine & 1.394 & 0.3071 & 1.501 & 0.282 \\
Raw + locality & 3.879 & 0.3038 & 3.327 & 0.2798 \\
Raw + KAHM & 2.332 & 0.3819 & 2.022 & 0.3359 \\
Seasonal naive & 1.723 & 0.2078 & 1.723 & 0.1725 \\
Holt & 2.631 & 0.345 & 2.632 & 0.3683 \\
Histogram boosting & 9.795 & 3.468 & 10.23 & 2.326 \\
MLP & 1.253 & 1.108 & 1.302 & 0.7098 \\
Always raw & 4.582 & 0.1846 & 2.984 & 0.1838 \\
Always raw + KAHM & 2.332 & 0.3819 & 2.022 & 0.3359 \\
Validation: matched pair & 2.342 & 0.5295 & 2.493 & 0.4758 \\
Validation: raw portfolio & 1.128 & 4.223 & 1.127 & 3.142 \\
Validation: KAHM family & 2.261 & 0.6594 & 2.021 & 0.5963 \\
Validation: full portfolio & 1.137 & 5.805 & 1.166 & 4.488 \\
Cheap probe & 2.297 & 0.4321 & 2.463 & 0.4022 \\
Metadata & 2.332 & 0.3828 & 2.022 & 0.3367 \\
Metadata + geometry & 2.364 & 0.4492 & 2.072 & 0.404 \\
Metadata + probe & 2.337 & 0.419 & 2.007 & 0.3846 \\
Metadata + probe + geometry & 2.348 & 0.4327 & 2.03 & 0.388 \\
\bottomrule
\end{longtable}
}

{\small\setlength{\tabcolsep}{4pt}
\begin{longtable}{lrrrr}
\caption{Confirmation capacity comparison. Each capacity changes specified training-window and model-size caps; it is not an equal-time budget. Values are equal-source weighted and retain all prescribed context--horizon settings and seeds.}\label{tab:capacity-confirmation}\\
\toprule
Method or policy & medium MSE & medium wall (s) & small MSE & small wall (s) \\
\midrule
\endfirsthead
\toprule
Method or policy & medium MSE & medium wall (s) & small MSE & small wall (s) \\
\midrule
\endhead
Raw ridge & 0.7635 & 0.1561 & 0.7961 & 0.2199 \\
Raw RBF & 0.7767 & 0.1719 & 0.8383 & 0.1657 \\
Affine local & 0.916 & 0.3088 & 0.9738 & 0.6917 \\
Locality local & 0.8908 & 0.3176 & 0.902 & 0.2621 \\
KAHM local & 0.9243 & 0.334 & 0.9466 & 0.3097 \\
Raw + affine & 0.8233 & 0.3859 & 0.8684 & 0.2899 \\
Raw + locality & 0.8363 & 0.2831 & 0.8561 & 0.275 \\
Raw + KAHM & 0.8288 & 0.3732 & 0.8551 & 0.4016 \\
Seasonal naive & 1.05 & 0.1621 & 1.05 & 0.1607 \\
Holt & 2.373 & 0.333 & 2.373 & 0.3569 \\
Histogram boosting & 0.9215 & 12.73 & 0.9896 & 5.086 \\
MLP & 0.8712 & 1.956 & 0.8859 & 0.9695 \\
Always raw & 0.7767 & 0.1719 & 0.8383 & 0.1657 \\
Always raw + KAHM & 0.8288 & 0.3732 & 0.8551 & 0.4016 \\
Validation: matched pair & 0.7847 & 0.4538 & 0.8347 & 0.5085 \\
Validation: raw portfolio & 0.8404 & 11.6 & 0.8874 & 5.074 \\
Validation: KAHM family & 0.8281 & 0.6116 & 0.8531 & 0.6369 \\
Validation: full portfolio & 0.8541 & 12.98 & 0.8987 & 6.918 \\
Cheap probe & 0.7943 & 0.3749 & 0.8343 & 0.3182 \\
Metadata & 0.8009 & 0.3187 & 0.8311 & 0.2764 \\
Metadata + geometry & 0.8009 & 0.3779 & 0.8311 & 0.3611 \\
Metadata + probe & 0.7765 & 0.2805 & 0.8377 & 0.3812 \\
Metadata + probe + geometry & 0.7765 & 0.2805 & 0.8377 & 0.3812 \\
\bottomrule
\end{longtable}
}

{\small\setlength{\tabcolsep}{4pt}
\begin{longtable}{lrrrr}
\caption{Estimated-benefit diagnostics with frozen hierarchical weights. These are descriptive log-benefit prediction errors and classification accuracies; they are not calibrated probabilities or interval coverage. No recalibration is fitted on confirmation data.}\label{tab:diagnostic-calibration}\\
\toprule
Diagnostic & Log RMSE & Log bias & Sign (\%) & 5\% class (\%) \\
\midrule
\endfirsthead
\toprule
Diagnostic & Log RMSE & Log bias & Sign (\%) & 5\% class (\%) \\
\midrule
\endhead
Development: Metadata & 1.035 & 0.01412 & 57 & 49.2 \\
Development: Metadata + geometry & 1.037 & 0.03491 & 47.7 & 43 \\
Development: Metadata + probe & 1.012 & 0.03064 & 52.3 & 53.9 \\
Development: Metadata + probe + geometry & 1.02 & 0.05091 & 50.8 & 43 \\
Confirmation: Metadata & 0.1643 & 0.05835 & 55.5 & 59.4 \\
Confirmation: Metadata + geometry & 0.1612 & 0.05269 & 55.5 & 75 \\
Confirmation: Metadata + probe & 0.1845 & 0.04069 & 58.6 & 73.4 \\
Confirmation: Metadata + probe + geometry & 0.1847 & 0.04196 & 58.6 & 71.9 \\
\bottomrule
\end{longtable}
}

{\small\setlength{\tabcolsep}{4pt}
\begin{longtable}{lrrrr}
\caption{One-off diagnostic development costs. Acquisition sums the union of both matched predictors' complete search/final nodes and each variant's metadata/probe/geometry dependencies within every development task. Summing standalone tasks can conservatively repeat preparation across seeds/tasks. Acquire and total denote wall time; CPU covers acquisition. Do not add the four rows: their developmental labels and some features can be shared. Ridge fitting is measured separately; no deployment-count amortization is assumed.}\label{tab:offline-costs}\\
\toprule
Variant & Acquire (s) & CPU (s) & Ridge (ms) & Total (s) \\
\midrule
\endfirsthead
\toprule
Variant & Acquire (s) & CPU (s) & Ridge (ms) & Total (s) \\
\midrule
\endhead
Metadata & 68.49 & 68.45 & 6.657 & 68.5 \\
Metadata + geometry & 78.6 & 78.56 & 10.04 & 78.61 \\
Metadata + probe & 80.83 & 80.79 & 9.035 & 80.84 \\
Metadata + probe + geometry & 80.83 & 80.79 & 6.146 & 80.84 \\
\bottomrule
\end{longtable}
}

{\small\setlength{\tabcolsep}{4pt}
\begin{longtable}{lrrrrrr}
\caption{Frozen development-only standardization and ridge coefficients, rounded for display. G denotes geometry and P the cheap probe. An absent feature is shown by a dash. Predictions use the unrounded values in the frozen diagnostic JSON and the fixed zero decision threshold.}\label{tab:diagnostic-coefficients}\\
\toprule
Feature & Mean & Scale & Meta & Meta+G & Meta+P & Meta+P+G \\
\midrule
\endfirsthead
\toprule
Feature & Mean & Scale & Meta & Meta+G & Meta+P & Meta+P+G \\
\midrule
\endhead
Intercept & -- & -- & 0.023966 & 0.023966 & 0.023966 & 0.023966 \\
log(1 + training windows) & 5.1417 & 0.66192 & 0.017096 & 0.013047 & -0.0099182 & -0.0089504 \\
log(1 + horizon/history) & 0.23442 & 0.019538 & -0.012299 & -0.013178 & -0.0043432 & -0.0045441 \\
Lag-one autocorrelation & 0.7189 & 0.24263 & 0.0069415 & 0.001093 & -0.014509 & -0.012806 \\
KAHM gate entropy & 0.074168 & 0.063288 & -- & -0.010312 & -- & -0.0026401 \\
log(1 + affine residual) & 0.18213 & 0.12986 & -- & -0.013232 & -- & 0.0063025 \\
Probe log-loss ratio & 0.25347 & 0.67921 & -- & -- & 0.14611 & 0.1467 \\
\bottomrule
\end{longtable}
}

{\small\setlength{\tabcolsep}{4pt}
\begin{longtable}{lrrrrr}
\caption{Recorded fitting behavior. Search and selected-refit warning events are separated. Counts describe task fits, which include dependent seeds and settings. Warned counts fits with recorded warnings; at cap counts selected refits reaching the recorded iteration limit. Full warning text is retained in the JSON report. No warning-triggered post hoc retuning is applied.}\label{tab:fit-warnings}\\
\toprule
Method & Fits & Warned & Search & Refit & At cap \\
\midrule
\endfirsthead
\toprule
Method & Fits & Warned & Search & Refit & At cap \\
\midrule
\endhead
Development: Raw ridge & 128 & 0 & 0 & 0 & 0 \\
Development: Raw RBF & 128 & 0 & 0 & 0 & 0 \\
Development: Affine local & 128 & 0 & 0 & 0 & 0 \\
Development: Locality local & 128 & 0 & 0 & 0 & 0 \\
Development: KAHM local & 128 & 0 & 0 & 0 & 0 \\
Development: Raw + affine & 128 & 0 & 0 & 0 & 0 \\
Development: Raw + locality & 128 & 0 & 0 & 0 & 0 \\
Development: Raw + KAHM & 128 & 0 & 0 & 0 & 0 \\
Development: Seasonal naive & 128 & 0 & 0 & 0 & 0 \\
Development: Holt & 128 & 0 & 0 & 0 & 0 \\
Development: Histogram boosting & 128 & 0 & 0 & 0 & 0 \\
Development: MLP & 128 & 128 & 232 & 105 & 105 \\
Confirmation: Raw ridge & 80 & 0 & 0 & 0 & 0 \\
Confirmation: Raw RBF & 80 & 0 & 0 & 0 & 0 \\
Confirmation: Affine local & 80 & 0 & 0 & 0 & 0 \\
Confirmation: Locality local & 80 & 0 & 0 & 0 & 0 \\
Confirmation: KAHM local & 80 & 0 & 0 & 0 & 0 \\
Confirmation: Raw + affine & 80 & 0 & 0 & 0 & 0 \\
Confirmation: Raw + locality & 80 & 0 & 0 & 0 & 0 \\
Confirmation: Raw + KAHM & 80 & 0 & 0 & 0 & 0 \\
Confirmation: Seasonal naive & 80 & 0 & 0 & 0 & 0 \\
Confirmation: Holt & 80 & 0 & 0 & 0 & 0 \\
Confirmation: Histogram boosting & 80 & 0 & 0 & 0 & 0 \\
Confirmation: MLP & 80 & 79 & 137 & 68 & 68 \\
\bottomrule
\end{longtable}
}

{\small\setlength{\tabcolsep}{4pt}
\begin{longtable}{lrrrrr}
\caption{Selection at numerical search-grid endpoints. Counts are descriptive; a two-point grid necessarily chooses an endpoint. These observations do not establish the benefit of a larger grid, and confirmation outcomes are not used to retune any method.}\label{tab:grid-boundaries}\\
\toprule
Method & Parameter & Choices & Lower & Upper & Declared grid \\
\midrule
\endfirsthead
\toprule
Method & Parameter & Choices & Lower & Upper & Declared grid \\
\midrule
\endhead
Development: Raw ridge & alpha & 128 & 84 & 4 & 0.01,1,100 \\
Development: Raw RBF & alpha & 128 & 35 & 4 & 0.01,1,100 \\
Development: Affine local & alpha & 128 & 45 & 8 & 0.01,1,100 \\
Development: Locality local & alpha & 128 & 35 & 9 & 0.01,1,100 \\
Development: KAHM local & alpha & 128 & 51 & 7 & 0.01,1,100 \\
Development: Raw + affine & alpha & 128 & 32 & 4 & 0.01,1,100 \\
Development: Raw + locality & alpha & 128 & 25 & 5 & 0.01,1,100 \\
Development: Raw + KAHM & alpha & 128 & 30 & 4 & 0.01,1,100 \\
Development: Histogram boosting & regularization & 128 & 78 & 50 & 0.01,1 \\
Development: MLP & regularization & 128 & 26 & 102 & 0.01,1 \\
Confirmation: Raw ridge & alpha & 80 & 20 & 56 & 0.01,1,100 \\
Confirmation: Raw RBF & alpha & 80 & 13 & 56 & 0.01,1,100 \\
Confirmation: Affine local & alpha & 80 & 9 & 58 & 0.01,1,100 \\
Confirmation: Locality local & alpha & 80 & 9 & 60 & 0.01,1,100 \\
Confirmation: KAHM local & alpha & 80 & 8 & 59 & 0.01,1,100 \\
Confirmation: Raw + affine & alpha & 80 & 12 & 59 & 0.01,1,100 \\
Confirmation: Raw + locality & alpha & 80 & 12 & 61 & 0.01,1,100 \\
Confirmation: Raw + KAHM & alpha & 80 & 12 & 60 & 0.01,1,100 \\
Confirmation: Histogram boosting & regularization & 80 & 34 & 46 & 0.01,1 \\
Confirmation: MLP & regularization & 80 & 5 & 75 & 0.01,1 \\
\bottomrule
\end{longtable}
}

